% Documento autonomo para Overleaf. Compilador: pdfLaTeX.
% No necesita archivos externos, BibTeX, minted ni shell-escape.
\documentclass[11pt,a4paper]{article}
\usepackage[utf8]{inputenc}
\usepackage[T1]{fontenc}
\usepackage{lmodern}
\usepackage[spanish,es-tabla,es-nodecimaldot]{babel}
\usepackage[margin=2.4cm,headheight=15pt]{geometry}
\usepackage{microtype}
\usepackage{amsmath}
\usepackage{xcolor}
\usepackage{booktabs,tabularx,longtable,array}
\usepackage{enumitem}
\usepackage{listings}
\usepackage{tikz}
\usetikzlibrary{arrows.meta,positioning,babel}
\usepackage{xurl}
\usepackage{hyperref}
\hypersetup{colorlinks=true,linkcolor=blue!45!black,urlcolor=blue!55!black,
  citecolor=blue!45!black,pdftitle={NAT Router Lite: funcionamiento y decisiones tecnicas},
  pdfauthor={Proyecto NAT Router Lite}}
\definecolor{codebg}{RGB}{245,247,250}
\lstset{basicstyle=\ttfamily\footnotesize,backgroundcolor=\color{codebg},
  frame=single,rulecolor=\color{black!15},breaklines=true,
  columns=fullflexible,keepspaces=true,showstringspaces=false,
  tabsize=4,aboveskip=1em,belowskip=1em}
\setlist{itemsep=3pt,topsep=5pt}
\setlength{\parindent}{0pt}
\setlength{\parskip}{6pt}
\setlength{\emergencystretch}{2em}
\renewcommand{\arraystretch}{1.18}
\newcolumntype{Y}{>{\raggedright\arraybackslash}X}
\newcommand{\code}[1]{\texttt{\detokenize{#1}}}
\newcommand{\decision}[1]{\par\textbf{Razón de diseño.} #1\par}
\title{\textbf{NAT Router Lite}\\[0.4em]
  \large Funcionamiento, arquitectura y decisiones técnicas}
\author{Documentación del proyecto}
\date{Base técnica: ESP-IDF 6.1\\ESP32 y ESP32-S3}

\begin{document}
\maketitle
\begin{abstract}
NAT Router Lite convierte un ESP32 o ESP32-S3 en un router Wi-Fi IPv4
con configuración mediante navegador. El dispositivo mantiene un punto
de acceso para sus clientes y una estación conectada al módem, guarda las
credenciales en NVS y comparte la salida mediante NAPT. Este documento
describe el código actual, sus interfaces y las razones técnicas de su
organización: eventos serializados, DNS local estable, confirmación HTTP,
diagnóstico de errores e indicador configurable. También delimita qué
comprueba el estado mostrado y qué funciones todavía no existen.
Las explicaciones de diseño se basan en el comportamiento observable del
código; no constituyen mediciones de rendimiento.
\end{abstract}
\tableofcontents
\clearpage

\section{Objetivo y recorrido de uso}
\label{sec:uso}

El objetivo es permitir que uno o varios clientes utilicen la conexión
Wi-Fi de un módem a través del ESP, sin recompilar cada vez que cambien
las credenciales del módem. El firmware crea la red \code{ESP32-NAT}, con
clave inicial \code{12345678}, y expone su aplicación en
\url{http://192.168.2.1}. El límite configurado es de tres clientes
asociados al AP. Tanto el ESP32 como el ESP32-S3 de este proyecto trabajan
con redes de 2,4 GHz.

El primer uso sigue esta secuencia:
\begin{enumerate}
  \item El usuario conecta el teléfono al AP del ESP. Si todavía no hay
  salida a Internet, debe mantener esa conexión para configurar el equipo.
  \item Abre la dirección HTTP anterior e introduce el SSID y la clave
  de la red del módem. El botón con un ojo permite mostrar u ocultar la clave.
  \item El servidor valida y guarda los datos. Una redirección conduce a
  una página que confirma la recepción y consulta el estado automáticamente.
  \item El ESP intenta asociarse, obtener una IPv4 y preparar NAPT y DNS.
  La página muestra espera, un diagnóstico o el estado preparado.
  \item Al reiniciar, se cargan las credenciales guardadas. Para cambiarlas
  se utiliza el enlace de la página; no es necesario borrar NVS.
\end{enumerate}

Guardar datos, conectarse a una red y disponer de Internet son hechos
distintos. El guardado puede ser correcto aunque la contraseña sea
incorrecta para ese módem. Una asociación correcta puede carecer de
IPv4 por un problema DHCP. Incluso con IPv4, NAT y DNS preparados, el
módem puede haber perdido su servicio de Internet. La aplicación informa
su estado local y las razones comunicadas por el controlador; no realiza
una prueba activa contra Google, YouTube u otro servicio externo.

\section{Topología y separación entre control y tráfico}
\label{sec:topologia}

\begin{figure}[htbp]
\centering
\begin{tikzpicture}[
  box/.style={draw=blue!45!black,rounded corners,fill=blue!4,
    minimum height=1.5cm,text width=3.5cm,align=center,font=\small},
  flow/.style={-{Latex[length=2mm]},thick},node distance=1.15cm]
  \node[box] (client) {Teléfono / PC\\Cliente del AP\\IPv4 privada por DHCP};
  \node[box,right=of client] (esp) {ESP32 / ESP32-S3\\AP: 192.168.2.1/24\\STA: IPv4 del módem};
  \node[box,right=of esp] (modem) {Módem / router\\DHCP y DNS de origen\\Salida a Internet};
  \draw[flow] (client) -- node[above,font=\small]{Wi-Fi} (esp);
  \draw[flow] (esp) -- node[above,font=\small]{Wi-Fi} (modem);
\end{tikzpicture}
\caption{Dos interfaces lógicas sobre una radio Wi-Fi. Las flechas representan el recorrido de salida.}
\end{figure}

\textbf{Plano de control.} El navegador habla con el servidor HTTP local,
que valida formularios, solicita almacenamiento y publica estado. Ese
intercambio termina en el ESP.

\textbf{Plano de datos.} El tráfico de las aplicaciones del teléfono
atraviesa lwIP y NAPT hacia la interfaz STA. Los paquetes de YouTube o
Google no se retransmiten mediante el servidor HTTP del proyecto.
El proxy DNS sí termina consultas UDP locales y las envía de nuevo al
servidor DNS seleccionado.

El modo \code{WIFI_MODE_APSTA} mantiene AP y STA simultáneamente. No
equivale a disponer de dos radios independientes: el canal de la red a
la que se asocia la STA tiene prioridad y el AP debe acompañarlo
\cite{wifi}. Por ello, el canal inicial 1 del AP no es una promesa de que
permanecerá en 1. El cambio de canal puede afectar transitoriamente a
clientes que no sigan adecuadamente el anuncio del AP.

\decision{APSTA permite conservar el acceso a la configuración mientras
la conexión al módem falla o se modifica. NAPT permite usar una subred
propia sin exigir rutas adicionales en el módem. A cambio, el resultado
es un router de capa IP con traducción de puertos, no un puente
transparente de capa 2. El rendimiento comparte los recursos de radio,
CPU y memoria del mismo microcontrolador.}

\section{Organización del código}
\label{sec:componentes}

\begin{tabularx}{\textwidth}{@{}p{2.7cm}YY@{}}
\toprule
Unidad & Responsabilidad & Dependencias propias \\
\midrule
\code{main} & Arranque y actualización del indicador & NVS, Wi-Fi, HTTP, LED \\
\code{nvs} & Leer y guardar credenciales & Ninguna \\
\code{wifi} & AP, STA, DHCP, eventos, NAPT, reconexión, diagnóstico y parser & NVS, proxy DNS \\
\code{server} & HTTP y recursos de la aplicación & Wi-Fi y su parser público \\
\code{dns_proxy} & Reenvío DNS UDP & Ninguna \\
\code{status_led} & Indicador GPIO o RGB mediante RMT & Ninguna \\
\bottomrule
\end{tabularx}

Cada componente declara en CMake las dependencias de ESP-IDF que utiliza.
Los encabezados públicos están en \code{include/};
\path{components/wifi/wifi_internal.h} permanece privado. Dentro de Wi-Fi,
\code{ap.c} configura AP/DHCP, \code{sta.c} configura la estación y
clasifica fallos, \code{wifi.c} coordina eventos y estado, y
\code{wifi_form.c} decodifica el formulario.

\path{main/main.c} inicializa LED, NVS, Wi-Fi y servidor en ese orden.
Después consulta el estado de conexión para actualizar el LED cada
500 ms. No contiene manejadores HTTP, parser, lógica DNS ni el flujo
completo de reconexión.

\decision{La división sigue responsabilidades y dependencias, de modo
que el servidor solicita acciones a Wi-Fi sin manipular sus variables
internas ni abrir NVS directamente. DNS conserva una tarea y un protocolo
propios; el LED solo recibe un booleano. Estas fronteras permiten probar
el parser o el proxy en host y cambiar una presentación sin trasladar
la lógica de red a la interfaz.}

El almacenamiento expone \code{router_nvs.h}, en lugar de un encabezado
local llamado \code{nvs.h}. Así se evita ocultar el encabezado homónimo
del SDK durante la resolución de inclusiones. El namespace y las claves
no se cambiaron al modularizar: se conserva compatibilidad con las
credenciales que ya estaban guardadas.

La API pública Wi-Fi incluye \code{router_wifi_start()},
\code{router_wifi_get_status()}, \code{router_wifi_is_connected()},
\code{router_wifi_store_credentials()} y
\code{router_wifi_apply_credentials()}. Guardar y aplicar son operaciones
separadas: la primera puede fallar en flash; la segunda puede fallar al
encolar el trabajo. La respuesta HTTP puede distinguir ambas situaciones.

\section{Inicialización y modelo de concurrencia}
\label{sec:eventos}

El arranque de Wi-Fi realiza las siguientes operaciones:
\begin{enumerate}
  \item Crea el temporizador de reconexión antes de habilitar los eventos.
  \item Inicializa ESP-NETIF, el bucle de eventos y sus manejadores.
  \item Inicializa el controlador y selecciona APSTA.
  \item Crea las interfaces AP y STA, configura DHCP y carga credenciales.
  \item Inicia el listener DNS en la dirección del AP.
  \item Ejecuta \code{esp_wifi_start()} y espera su retorno correcto.
  \item Publica \code{ROUTER_CONNECT}; el manejador de eventos inicia
  entonces la conexión STA.
\end{enumerate}

El listener DNS existe antes de que un teléfono pueda obtener su primera
concesión DHCP. Por otra parte, la primera conexión se solicita después
de que termine \code{esp_wifi_start()}. En las pruebas originales del S3,
conectar directamente al recibir \code{WIFI_EVENT_STA_START} llegó a
producir \code{ESP_ERR_WIFI_NOT_STARTED}. La secuencia actual evita esa
carrera de arranque en este proyecto; no significa que el patrón de
conectar desde ese evento esté prohibido en todos los programas de IDF.

\subsection{Una tarea coordina las operaciones del controlador}

El callback del temporizador publica un evento \code{ROUTER_CONNECT};
no ejecuta directamente \code{esp_wifi_connect()}. El manejador HTTP
publica \code{ROUTER_UPDATE_CREDENTIALS} mediante la API pública. El
bucle de eventos procesa esas solicitudes y las notificaciones Wi-Fi/IP.

\decision{Serializar las operaciones de conexión, desconexión y cambio
de configuración reduce carreras entre la tarea HTTP, la tarea del
temporizador y el controlador. El callback del temporizador queda corto
y el servidor puede contestar sin esperar toda la asociación. Esto no
convierte el sistema en un programa de una sola tarea: HTTP, DNS,
temporizadores, Wi-Fi y lwIP siguen teniendo ejecución concurrente.}

Los indicadores compartidos con HTTP utilizan atómicos C11:
conectado, credenciales guardadas, NAPT, conflicto y última razón de
desconexión. Los campos se leen individualmente; la estructura resultante
no es una fotografía transaccional de todos ellos. Durante una transición
una consulta puede observar una combinación pasajera que se estabiliza
en el siguiente sondeo. Las banderas privadas que gobiernan la aplicación
de credenciales se usan desde la tarea de eventos.

\subsection{Reconexión y actualización de credenciales}

Tras una desconexión se marca la salida como no disponible: se limpia
el estado de conexión, se deshabilita NAPT y se invalida el DNS de origen.
Un temporizador de una sola ejecución programa un nuevo intento después
de tres segundos. A esto se añaden el escaneo de canales y los reintentos
internos del controlador, configurados con \code{failure_retry_cnt = 2}.
Tres segundos es el intervalo programado entre intentos de la aplicación,
no la duración total máxima de una conexión.

Al recibir nuevas credenciales, el flujo es:
\begin{enumerate}
  \item Confirmar el guardado en NVS y limpiar el diagnóstico anterior.
  \item Encolar la actualización, marcar la salida sin servicio y detener
  el temporizador que pudiera estar pendiente.
  \item Solicitar desconexión STA, recuperar los datos guardados y aplicar
  la nueva configuración. Si todavía no se puede, mantener la actualización
  pendiente para un intento posterior.
  \item Programar la conexión con la configuración actualizada.
  \item Al obtener IPv4, preparar la ruta, NAPT y el DNS de origen.
\end{enumerate}

La desconexión provocada localmente por cambiar claves no debe mostrarse
como si el módem hubiera rechazado la nueva contraseña. El código reconoce
las razones de salida local esperadas y evita convertirlas en un fallo
de autenticación. Un fallo real posterior sí queda registrado.

\code{WIFI_EVENT_STA_CONNECTED} indica asociación, pero el booleano
\code{connected} se activa al recibir \code{IP_EVENT_STA_GOT_IP}.
\code{STA_STOP}, \code{STA_DISCONNECTED} y \code{STA_LOST_IP} invalidan
el estado de salida correspondiente. Si aparece una IP mientras hay un
cambio de credenciales pendiente, no se presenta como la conexión final
de esa actualización.

\section{Credenciales: validación y persistencia}
\label{sec:credenciales}

\subsection{Interpretación correcta del formulario}

El navegador envía \code{ssid} y \code{password} mediante un POST con
codificación de formulario URL. El servidor acumula todos los fragmentos
del cuerpo: una lectura de socket no garantiza recibir el mensaje entero.
Se rechaza un cuerpo vacío, de 320 bytes o más, con NUL embebido o con
campos inválidos. El búfer de 320 bytes contempla la expansión de cada
byte a tres caracteres \texttt{\%xx} y la sintaxis de ambos campos.

El parser separa primero los pares por \texttt{\&} y cada clave de su valor
por \code{=}; después decodifica. Un separador codificado dentro de una
contraseña debe seguir siendo parte de su valor. Por ejemplo:
\begin{lstlisting}
ssid=Mi+red&password=abc%26def%3D123
# SSID decodificado: Mi red
# Clave decodificada: abc&def=123
\end{lstlisting}
El signo \code{+} significa espacio; \texttt{\%2B} representa un signo más
literal. Las secuencias incompletas, no hexadecimales o que produzcan
\texttt{\%00} son inválidas. Los dos campos requeridos deben aparecer una
sola vez. Las claves desconocidas bien formadas se ignoran; no sustituyen
a los campos obligatorios.

Se aceptan SSID de 1 a 32 bytes y claves de 8 a 63 bytes, o una PSK de
64 caracteres hexadecimales. Los búferes de salida tienen 33 y 65 bytes
para incluir el terminador NUL. Son límites en bytes, no en caracteres
Unicode: un carácter acentuado puede ocupar más de un byte UTF-8.
Los atributos HTML ayudan al usuario, pero la validación definitiva se
realiza en C. No hay un modo de formulario para redes abiertas.

\decision{Validar en el servidor evita depender de un navegador concreto
o de sus límites de formulario. Separar antes de decodificar conserva
contraseñas que contienen \texttt{\&}, \code{=} o espacios. El parser es C
independiente de ESP-IDF y no modifica sus salidas si falla, por lo que
puede probarse con casos de borde en un ordenador.}

\subsection{Qué se guarda y cuándo se conserva}

El componente NVS utiliza la partición NVS predeterminada, el namespace
\code{wifi} y las claves de cadena \code{ssid} y \code{password}.
Abre en lectura/escritura, escribe ambas cadenas y llama a
\code{nvs_commit()}; solo informa éxito si todas esas operaciones
devuelven \code{ESP_OK}. Al cargar usa lectura solamente y comprueba
que las cadenas caben y cumplen las longitudes mínimas esperadas.
NVS es el almacenamiento de pares clave--valor en flash de ESP-IDF
\cite{nvs}.

El arranque normal no borra credenciales. La inicialización recupera
automáticamente la partición mediante borrado únicamente si recibe
\code{ESP_ERR_NVS_NO_FREE_PAGES} o
\code{ESP_ERR_NVS_NEW_VERSION_FOUND}. Otros fallos se propagan al arranque.
Si no hay datos válidos, la STA usa los marcadores \code{SSID}/\code{PASS}
y el AP permanece disponible para recibir los datos reales. Es posible
ver avisos de red no encontrada antes de configurar el equipo.

Guardar dos claves y hacer commit no añade por sí mismo una transacción
de aplicación que garantice una pareja indivisible ante cualquier corte
de energía o error intermedio. El código actual no utiliza un registro
versionado ni un único blob que agrupe SSID y clave. Si se necesitara esa
garantía adicional, habría que diseñar y validar ese formato de persistencia.

\decision{Dos cadenas pequeñas encajan con NVS y conservan compatibilidad
con los datos de la versión anterior. No se escribe flash en cada sondeo
de estado ni en cada actualización del LED: se guarda cuando se reciben
credenciales. La información sensible no se incorpora al JSON de estado
ni a la página de confirmación.}

\section{IPv4, DHCP y NAPT}
\label{sec:napt}

El AP tiene dirección y gateway \code{192.168.2.1}, con máscara
\code{255.255.255.0}. Su servidor DHCP entrega configuración a los
clientes. La interfaz STA utiliza el DHCP de la red del módem para
obtener su dirección, máscara, gateway y DNS.

Después de obtener IPv4, el firmware comprueba si las redes AP y STA
se solapan. Compara las direcciones de ambas interfaces bajo las dos
máscaras. Por ejemplo, si el módem también utiliza \code{192.168.2.0/24},
la salida no puede distinguir correctamente esa red del segmento de
clientes del ESP. También puede existir solapamiento con una máscara
más amplia aunque los números de subred parezcan distintos.

Cuando detecta conflicto, conserva el dato de que la STA obtuvo IP,
activa \code{subnet_conflict}, deja NAPT deshabilitado y retira el DNS de
origen. No cambia automáticamente la IP del AP. Resolverlo requiere
elegir subredes diferentes y ajustar la configuración de dirección en el
código; el filtro de clientes DNS usa igualmente una máscara /24.

Si las redes son compatibles, se selecciona la STA como interfaz de
salida predeterminada y se habilita NAPT sobre el AP mediante
\code{esp_netif_napt_enable()} \cite{netif}. Los defaults comunes activan:
\begin{lstlisting}
CONFIG_LWIP_IP_FORWARD=y
CONFIG_LWIP_IPV4_NAPT=y
CONFIG_LWIP_IPV4_NAPT_PORTMAP=y
\end{lstlisting}

NAPT mantiene traducciones de direcciones y puertos para los flujos de
salida de los clientes. El módem ve ese tráfico a través de la dirección
STA del ESP; las respuestas se devuelven al cliente que originó el flujo.
No se necesita abrir un socket BSD de la aplicación por cada conexión
reenviada: las traducciones pertenecen a la pila de red. Su tabla y sus
búferes sí consumen recursos de lwIP.

\decision{Habilitar NAPT únicamente con IPv4 útil evita presentar como
disponible una salida que todavía no tiene ruta. Deshabilitarlo al perder
la STA hace coherente el estado publicado. Detectar la superposición
evita atribuir a una contraseña o al DNS un problema de direccionamiento.}

El soporte de port mapping está compilado, pero la aplicación no crea
reglas de publicación de puertos ni proporciona una interfaz para ellas.
Tampoco implementa encaminamiento IPv6. Consultar un registro DNS AAAA
puede producir una respuesta válida sin que exista conectividad IPv6
para usarla.

\section{Proxy DNS: la razón del DNS local estable}
\label{sec:dns}

\subsection{Por qué cambiar el DNS de DHCP no bastaba}

Un teléfono puede asociarse al AP y obtener su concesión DHCP antes de
que el ESP se conecte al módem. La concesión incluye un servidor DNS que
el teléfono conserva. Modificar después la configuración del servidor
DHCP no obliga al teléfono a solicitar una concesión inmediatamente.
Si la dirección recibida no ofrece un DNS funcional, los dominios fallan
aunque NAPT ya permita tráfico por dirección IP.

La implementación mantiene una dirección estable: DHCP anuncia desde
el arranque \code{192.168.2.1} como DNS, y el proxy escucha allí en UDP53
antes de iniciar la radio. Cuando la STA obtiene conexión, solo cambia
el servidor al que el proxy reenvía consultas. No se reinicia DHCP para
intentar sustituir el DNS en concesiones existentes.

\decision{El contrato con el cliente permanece constante durante el
aprovisionamiento. Esta separación entre DNS anunciado y DNS de origen
resuelve el caso en que el teléfono solo podía navegar después de apagar
y encender su Wi-Fi. El cliente conserva gateway y DNS, mientras el ESP
actualiza su salida internamente.}

\subsection{Selección del servidor de origen}

Al obtener IPv4 se consultan los DNS \code{MAIN}, \code{BACKUP} y
\code{FALLBACK} de la STA, en ese orden. Se elige el primero que sea IPv4
y pase las comprobaciones del código: no ser la IP del AP, no pertenecer
al primer octeto 0 ni 127, y estar por debajo del rango multicast.
Si ninguno cumple, se utiliza \code{1.1.1.1} y se registra un aviso.

La selección es una búsqueda de configuración válida, no una comprobación
de disponibilidad. El proxy no cambia automáticamente al segundo DNS
cuando una consulta expira; tampoco prueba que la red permita alcanzar
el DNS público de reserva. Una red con DNS interno o restricciones de
salida puede necesitar ajustes.

\subsection{Recorrido de una consulta}

El proxy usa dos sockets UDP: uno enlazado a la IP del AP y puerto 53,
y otro conectado al DNS de origen. Una tarea FreeRTOS con pila de
4096 bytes y prioridad 5 atiende las operaciones mediante \code{select()}.

\begin{enumerate}
  \item Recibe el datagrama y comprueba su procedencia del segmento AP /24.
  Valida la cabecera, que sea una consulta ordinaria y que contenga una
  sola pregunta con nombre y campos dentro de los límites.
  \item Reserva una de las seis entradas pendientes. Guarda el cliente,
  su identificador original, la pregunta necesaria para responder con error
  y el instante de inicio.
  \item Reemplaza el identificador DNS por uno propio, único entre las
  entradas activas. El contador inicial parte de \code{esp_random()}.
  Reenvía el mensaje manteniendo el resto de su contenido.
  \item Recibe desde el socket asociado al DNS seleccionado y comprueba
  que sea una respuesta para una transacción pendiente. La estructura de
  cabecera, ID y códigos de respuesta sigue el formato DNS \cite{dnsrfc}.
  \item Restaura el identificador original y entrega el mensaje al cliente.
\end{enumerate}

Dos teléfonos pueden usar el mismo identificador DNS. La reescritura y
la tabla de pendientes evitan confundirlos, incluso si las respuestas
llegan en orden inverso. El socket conectado limita la recepción al par
remoto configurado; la coincidencia de identificador es una comprobación
de correlación y no sustituye autenticación criptográfica de DNS.

El socket de salida no se enlaza a la IP del AP: utiliza la ruta
predeterminada STA. Enlazarlo a la dirección de la red de clientes
introduciría una elección de origen inadecuada para el camino al módem.
En UDP, \code{connect()} selecciona un par; no realiza un intercambio
que demuestre que el servidor responde.

\subsection{Límites, errores y cambios de conexión}

\begin{tabularx}{\textwidth}{@{}p{4.5cm}Y@{}}
\toprule
Situación & Tratamiento actual \\
\midrule
Consulta inválida & Respuesta \code{FORMERR} cuando puede construirse \\
Sin DNS de origen, tabla llena o fallo de envío & \code{SERVFAIL} \\
Sin respuesta en dos segundos & Expiración y \code{SERVFAIL} al cliente \\
Consulta mayor de 4096 bytes & \code{FORMERR} \\
Respuesta mayor de 4096 bytes & Indicador \code{TC}; sin fallback TCP local \\
Cambio de DNS o pérdida de STA & Invalida disponibilidad, falla pendientes y recrea el socket de origen \\
Respuesta con ID desconocido & Se descarta \\
\bottomrule
\end{tabularx}

Se usan hasta seis pendientes y un búfer de 4097 bytes para detectar
mensajes que exceden el máximo aceptado de 4096. El proxy admite preguntas
ordinarias con nombre sin compresión; no pretende interpretar todas las
variantes posibles del protocolo. Los datos adicionales aceptados se
reenvían sin convertir cada tipo de registro a una estructura propia.
Por eso el transporte puede llevar preguntas A, AAAA o HTTPS sin que
el firmware tenga que conocer el significado de cada respuesta.

\decision{\code{SERVFAIL} representa indisponibilidad o fallo temporal;
\code{NXDOMAIN} afirmaría que el dominio no existe y podría dejar una
respuesta negativa en caché. Los límites fijos acotan memoria y tiempo de
retención en un microcontrolador. El precio es rechazar carga superior
a seis pendientes y no cubrir todos los tamaños o transportes de DNS.}

Marcar una respuesta como truncada evita entregar silenciosamente un
mensaje incompleto. Sin embargo, un cliente que intente continuar por
TCP en el puerto 53 no encontrará un servidor TCP en el proxy. Es una limitación real
para respuestas grandes, no una implementación completa de fallback.
No hay caché DNS, DoH, DoT ni validación DNSSEC propia. Los requisitos
generales del transporte TCP están descritos en RFC 7766 \cite{dnstcp}.

Los cambios de servidor usan una dirección y un número de generación
atómicos. \code{dns_ready} exige que el proxy esté iniciado y que el socket
de origen corresponda a la generación vigente. Así no se conserva un
estado preparado asociado al DNS de una conexión anterior. La tarea
revisa actividad con un intervalo de selección de 100 ms y, cuando falla
la creación del socket de origen, vuelve a intentarla aproximadamente
cada segundo. Ninguno de esos indicadores demuestra por sí solo salida
a Internet.

\section{Servidor HTTP y aplicación del navegador}
\label{sec:http}

\subsection{Recursos y contrato de las rutas}

Las páginas \path{app/main.html} y \path{app/saved.html} se incorporan
como texto al binario mediante \code{EMBED_TXTFILES}; el favicon se añade
como binario con \code{EMBED_FILES}. El componente servidor usa los
símbolos del enlazador para transmitirlos. En los HTML descuenta el NUL
que añade la incorporación de texto; para el icono transmite el tamaño
binario completo.

\decision{La aplicación es pequeña y estática. Embebiendo sus recursos
se evita montar un sistema de archivos, gestionar una partición web o
distribuir archivos por separado. Los recursos y su API viajan con la
misma versión del firmware. Cambiar HTML exige recompilar y programar
esa versión; no existe un cargador de páginas independiente.}

\begin{longtable}{@{}p{3.3cm}p{4.3cm}p{7.3cm}@{}}
\toprule
Petición & Respuesta normal & Comportamiento relevante \\
\midrule
\endhead
\code{GET /} & HTML del formulario & Tipo UTF-8, \code{no-cache}, \code{nosniff} \\
\code{POST /} & \code{303} a \code{/saved} & Valida, guarda y encola aplicación de las credenciales \\
\code{GET /saved} & HTML de confirmación & \code{no-store}; redirige a \code{/} si no hay credenciales guardadas \\
\code{GET /api/status} & JSON de estado & \code{no-store}; no incluye SSID ni contraseña \\
\code{GET /favicon.ico} & Icono & Caché de un año, \code{immutable}, \code{nosniff} \\
\bottomrule
\end{longtable}

Un POST inválido produce 400; un timeout al recibir el cuerpo, 408.
Si no se puede guardar, se responde 500. Si se guardó pero falló la
publicación del evento de aplicación, también se devuelve 500, con un
mensaje que distingue que los datos quedaron persistidos y pueden
aplicarse al reiniciar. La confirmación 303 se envía cuando ambas etapas
inmediatas tuvieron éxito; no espera a que el módem acepte la contraseña.

\decision{El patrón POST--redirección--GET permite recargar la página de
confirmación sin volver a enviar las credenciales. Mostrar primero el
guardado y luego el progreso evita dejar al usuario sin respuesta durante
el escaneo, DHCP y los reintentos. La autenticación Wi-Fi se informa de
forma asíncrona mediante la API de estado.}

\subsection{Estado publicado y significado de cada campo}

Una respuesta preparada puede ser:
\begin{lstlisting}
{
  "credentials_saved": true,
  "connected": true,
  "nat_enabled": true,
  "dns_ready": true,
  "subnet_conflict": false,
  "wifi_reason": 0,
  "connection_error": "none"
}
\end{lstlisting}

\begin{tabularx}{\textwidth}{@{}p{4.1cm}Y@{}}
\toprule
Campo & Interpretación \\
\midrule
\code{credentials_saved} & Hay credenciales cargadas válidas o el guardado actual terminó correctamente; no certifica su aceptación por el módem. \\
\code{connected} & La STA recibió IPv4 y todavía no se invalidó ese estado. \\
\code{nat_enabled} & La activación de NAPT devolvió éxito. \\
\code{dns_ready} & Proxy iniciado y socket UDP de origen preparado para la generación actual. \\
\code{subnet_conflict} & Se detectó superposición de las redes AP y STA. \\
\code{wifi_reason} & Última razón numérica Wi-Fi registrada, o cero sin diagnóstico vigente. \\
\code{connection_error} & Categoría estable derivada de esa razón para la interfaz. \\
\bottomrule
\end{tabularx}

La condición local de listo combina conexión, NAPT, DNS y ausencia de
conflicto. No comprueba resolución efectiva en cada sondeo ni acceso a
un sitio externo. Tampoco hay un watchdog de conectividad global que
reinicie el dispositivo si el módem conserva Wi-Fi pero pierde Internet.

\subsection{Sondeo automático y presentación}

La página de confirmación consulta la API cada dos segundos durante la
espera o el error, y cada diez segundos cuando está preparada. Cada
petición tiene un límite de cinco segundos. Esta periodicidad corresponde
al navegador que mantiene abierta la página; no es una tarea de
monitorización adicional que siga ejecutándose dentro del ESP después
de cerrar la pestaña.

El indicador de carga aparece en espera y desaparece en estado listo,
error o inactivo. No existe ya un botón de actualización manual.
Aunque desaparezca la animación al estar preparado, el sondeo más lento
permite reflejar una pérdida posterior. Un fallo de petición HTTP no
demuestra por sí solo que la contraseña sea incorrecta. Pasados treinta
segundos sin un diagnóstico específico se muestra una espera prolongada.
La preferencia de movimiento reducido del sistema desactiva la animación.

El botón del ojo utiliza SVG en línea, cambia el tipo del campo entre
\code{password} y \code{text} y actualiza sus atributos accesibles.
Es un botón de tipo \code{button}, de modo que usarlo no envía el
formulario. Se desactivan autocorrección, capitalización automática y
revisión ortográfica para reducir modificaciones involuntarias del dato.

\decision{Sondear un pequeño JSON mantiene sencilla la interacción con
un servidor embebido y permite verificar el comportamiento en host.
No requiere WebSocket ni una conexión de eventos permanente. El intervalo
más largo al estar listo reduce peticiones; separar error de espera evita
una animación indefinida que esconda el diagnóstico.}

\section{Diagnóstico de asociación y autenticación}
\label{sec:diagnostico}

El controlador informa razones de desconexión; \code{sta.c} las agrupa
en categorías que la página transforma en mensajes comprensibles. Los
códigos representativos siguientes no agotan todos los motivos de cada
grupo; la tabla completa vigente está en \code{router_wifi_error_for_reason()}.

\begin{longtable}{@{}p{4.8cm}p{2.1cm}p{8cm}@{}}
\toprule
Categoría & Ejemplos & Interpretación práctica \\
\midrule
\endhead
\code{none} & 0 & Sin error Wi-Fi registrado en este momento. \\
\code{network_not_found} & 201 & No se encontró el AP solicitado: revisar SSID, alcance y disponibilidad de la red. \\
\code{handshake_timeout} & 15, 204 & No terminó el intercambio de claves; una contraseña incorrecta es una causa posible. \\
\code{authentication_failed} & 202, 23 & Rechazo o fallo de autenticación; revisar clave y seguridad del módem. \\
\code{authentication_timeout} & 2 & Expiró la fase de autenticación. \\
\code{security_mismatch} & 210, 211 & Seguridad o cifrado incompatibles con la configuración aceptada. \\
\code{signal_too_weak} & 212 & No se encontró una red que cumpla el umbral de señal indicado por el controlador. \\
\code{association_failed} & 203 & Falló la asociación al AP. \\
\code{modem_busy} & 5 & El AP rechazó más estaciones asociadas. \\
\code{connection_lost} & 200, 4 & Pérdida de beacons, inactividad u otra salida de la red. \\
\code{connection_failed} & Otros & Fallo que no pertenece a una categoría más específica. \\
\bottomrule
\end{longtable}

No se puede asegurar en todos los equipos que un timeout sea
exclusivamente una contraseña incorrecta. Interferencias, pérdida de
tramas o ajustes del módem pueden impedir completar el mismo intercambio.
La página orienta hacia una causa probable y conserva el código de
diagnóstico; no inventa una certeza que el controlador no proporciona.
La ausencia de un mensaje exacto no es una política deliberada de
ocultar el error por seguridad.

La última razón real permanece durante los reintentos automáticos, en
lugar de borrarse cada tres segundos y hacer que el mensaje desaparezca.
Se limpia al guardar nuevas credenciales o al alcanzar una asociación
correcta. Las salidas locales esperadas al cambiar configuración se
tratan por separado, como se explicó en la sección~\ref{sec:eventos}.

\decision{Una clasificación estable desacopla la interfaz de decenas de
constantes del SDK. Conservar razón numérica y categoría permite una
explicación breve al usuario y un dato más preciso para revisar el log.
La persistencia temporal del diagnóstico mejora su visibilidad sin
detener los intentos de recuperación.}

\section{Indicador y opciones específicas del hardware}
\label{sec:led}

El indicador es seleccionable en el menú \textbf{NAT Router}: sin LED,
LED GPIO o RGB direccionable. El ESP32 utiliza por defecto GPIO2 activo
bajo. El perfil S3 utiliza RGB en GPIO48, con brillo inicial 8. En
ESP32-S3-DevKitC-1 v1.1 el RGB se trasladó a GPIO38; el pin debe ajustarse
a la revisión física \cite{devkit}.

Sin IPv4 STA, el GPIO parpadea y el RGB parpadea en ámbar. Con IPv4, el
GPIO permanece encendido y el RGB pasa a verde. La actualización se
realiza cada 500 ms desde \code{app_main()}: un ciclo encendido/apagado
completo en desconexión dura aproximadamente un segundo. El LED refleja
\code{router_wifi_is_connected()}, no las tres condiciones de conexión,
NAPT y DNS juntas. Puede verse verde aunque el módem no tenga Internet o
exista conflicto de subred.

\subsection{Por qué se utiliza RMT para el RGB}

El LED direccionable requiere pulsos precisos. El código utiliza el
controlador RMT de transmisión de ESP-IDF y un copy encoder
\cite{rmt}; no depende de temporización manual de GPIO dentro de una
tarea ni de un componente externo de tiras LED. Configura una resolución
de 10 MHz: cada tick dura 0,1 microsegundos.

Se construyen 24 símbolos para los bits de color en orden GRB y un símbolo
adicional de reset. Para un cero se generan 3 ticks altos y 9 bajos;
para un uno, 7 altos y 6 bajos. El reset ocupa 300 microsegundos en bajo.
Los 25 símbolos caben en los 48 símbolos de memoria de un canal del S3,
por lo que no se necesita DMA ni rellenar una trama más larga por bloques.
La cola de transmisión tiene profundidad 1.

La trama se almacena estáticamente para mantener válido el búfer durante
la transmisión, incluso si se agota la espera. Si el GPIO no es válido,
falla la reserva del canal o se produce un error de actualización, se
registra el problema y se desactiva el indicador.

\decision{El LED es una ayuda de observación y su fallo no debe impedir
el router. RMT proporciona generación de pulsos por hardware mientras
Wi-Fi y otras tareas están activas. Limitarse a una trama completa pequeña
reduce la complejidad del codificador y evita exigir DMA solo para un LED.}

\subsection{FTM y consola USB}

El menú incluye un responder FTM opcional, condicionado por la capacidad
del target y las opciones del SDK. Permite responder a clientes Wi-Fi RTT
compatibles; está desactivado en los perfiles base. No añade una interfaz
de medición de distancia ni aumenta la velocidad del router.

El perfil S3 conserva la consola secundaria USB Serial/JTAG. Es una
opción de observación y depuración del hardware, no una función web ni
un intérprete de comandos remoto de la aplicación. PSRAM no se activa
por suponer que todos los módulos S3 la incluyen: debe verificarse el
módulo y su interfaz física.

\section{Perfiles de ESP-IDF y compilación reproducible}
\label{sec:configuracion}

\subsection{Chip, placa, defaults y configuración generada}

ESP32 y ESP32-S3 requieren seleccionar targets distintos. Además, un
mismo chip puede montarse en placas con flash, PSRAM o LED diferentes.
Los perfiles incluidos son preferencias de las placas del proyecto,
no una descripción universal de todos los módulos de cada familia.

\begin{tabularx}{\textwidth}{@{}p{5.4cm}Y@{}}
\toprule
Archivo & Papel \\
\midrule
\path{sdkconfig.defaults} & Forwarding IPv4, NAPT y soporte de port mapping comunes. \\
\path{sdkconfig.defaults.esp32} & DIO/40 MHz, flash declarada de 2 MB, GPIO2 activo bajo. \\
\path{sdkconfig.defaults.esp32s3} & DIO/80 MHz, flash de 16 MB, RGB48 y opciones específicas del S3 usado. \\
\path{sdkconfig} & Configuración activa completa generada por ESP-IDF. \\
\path{sdkconfig.old} & Respaldo automático previo; no es un perfil permanente de un chip. \\
\bottomrule
\end{tabularx}

ESP-IDF carga los defaults comunes y después el sufijo correspondiente a
\code{IDF_TARGET} \cite{build}. Una selección explícita ya guardada en
\code{sdkconfig} tiene prioridad sobre defaults. Por eso editar un perfil
no fuerza necesariamente todos sus valores en una configuración ya
existente. En IDF 6.1 las opciones con procedencia \texttt{\# default:}
pueden adoptar nuevos defaults durante la reconfiguración.

El campo \code{CONFIG_IDF_TARGET} identifica el chip de cada configuración
generada. Al organizar estos perfiles, el archivo raíz activo correspondía
a \code{esp32} y el respaldo a \code{esp32s3}; esos contenidos pueden
cambiar al volver a configurar. El significado estable de \code{.old}
es respaldo, no S3. No deben copiarse configuraciones completas entre
chips ni editarse manualmente las banderas de capacidades \code{SOC_*}.

El perfil S3 conserva la consola USB secundaria,
\code{CONFIG_FREERTOS_IN_IRAM}, colores del bootloader y colores del log.
El tamaño de 16 MB se apoyó en la detección de flash del equipo probado.
Los 2 MB del perfil ESP32 son el tamaño declarado en su configuración;
no certifican la capacidad física de cualquier placa ESP32.
40/80 MHz indican reloj de la memoria flash, no la banda de radio Wi-Fi.
El tamaño de flash tampoco aumenta por sí mismo la partición de aplicación:
el esquema de particiones debe revisarse si se requiere una aplicación mayor.

\decision{Los sufijos estándar de ESP-IDF hacen explícitas las diferencias
de placa sin convertir un respaldo generado en una segunda configuración
oficial. Mantener lo común en el archivo base reduce divergencias de
forwarding o NAPT entre chips. Conservar el GPIO configurable evita ligar
la lógica de red a una revisión concreta del DevKit.}

\subsection{Comandos de trabajo}

Desde una terminal con el entorno ESP-IDF 6.1 activado, el flujo habitual
consiste en seleccionar un único target, configurar, compilar y programar:
\begin{lstlisting}
# Elegir esp32 O esp32s3, segun el chip real.
idf.py set-target esp32s3
idf.py menuconfig
idf.py build
idf.py -p COM3 flash monitor
\end{lstlisting}
\code{COM3} debe reemplazarse por el puerto real. Se sale del monitor
con \code{Ctrl+]}. \code{set-target} limpia el directorio de compilación
seleccionado y regenera la configuración; no es un nombre alternativo
para elegir la placa al programar un binario ya existente. Si el entorno
define \code{IDF_TARGET}, debe coincidir con el target solicitado.

Para trabajar con ambos chips sin compartir configuración generada ni
caché, pueden utilizarse carpetas de compilación separadas:
\begin{lstlisting}
# PowerShell, desde la raiz del proyecto.
Remove-Item Env:IDF_TARGET -ErrorAction SilentlyContinue

idf.py -B build/esp32 -D IDF_TARGET=esp32 -D SDKCONFIG=build/esp32/sdkconfig build
idf.py -B build/esp32s3 -D IDF_TARGET=esp32s3 -D SDKCONFIG=build/esp32s3/sdkconfig build

# Programar especificamente el resultado S3.
idf.py -B build/esp32s3 -D IDF_TARGET=esp32s3 -D SDKCONFIG=build/esp32s3/sdkconfig -p COM3 flash monitor
\end{lstlisting}
Las mismas opciones deben acompañar a \code{menuconfig}, \code{build},
\code{flash} y \code{monitor}. Si una carpeta tiene caché de otro target
o de otra ruta \code{SDKCONFIG}, debe elegirse una carpeta nueva o
regenerarla de forma consciente.

\code{idf.py save-defconfig} escribe por defecto en el archivo raíz
\code{sdkconfig.defaults}; no distribuye automáticamente las opciones
entre los dos sufijos. Antes de conservar sus resultados deben separarse
preferencias comunes y de placa. La guía \path{CONFIGURACION.md} contiene
los comandos de mantenimiento y las comprobaciones de los perfiles.

\section{Presupuesto de recursos y tratamiento de fallos}
\label{sec:recursos}

Con la configuración revisada, lwIP tiene diez sockets disponibles.
El servidor fija cuatro sesiones HTTP abiertas como máximo y admite
purga de la sesión menos recientemente usada. Según la reserva interna
del servidor de IDF, se presupuestan tres sockets adicionales de HTTP;
el proxy usa dos. El total presupuestado es por tanto
\(4 + 3 + 2 = 9\), dentro del límite de diez.

\decision{Dejar el máximo HTTP predeterminado sin tener en cuenta DNS
podría agotar los sockets al iniciar ambos servicios. El límite de cuatro
sesiones y la purga LRU equilibran conexiones de navegador y el servicio
DNS. No representa cuatro clientes Wi-Fi: una sola pestaña puede utilizar
más de una conexión HTTP y el AP mantiene su límite independiente de tres
estaciones.}

Las seis entradas DNS y el búfer de paquetes son estáticos y pertenecen
a la tarea DNS. Evitan alojar el paquete de unos 4 KB en su pila de 4 KB.
Entre paquete y tabla pendiente se reserva del orden de 6 KB estáticos,
además de la pila de tarea. Esto es una estimación de esas estructuras,
no el consumo total de RAM: faltan TCB, sockets, pbufs, controlador
Wi-Fi, servidor, otras tareas y alineación. La disponibilidad real se
comprueba en el mapa de enlace y las métricas de heap.

El servidor reserva ocho posiciones de URI y registra cinco rutas.
Los recursos HTML ocupan flash del firmware y se sirven desde sus
símbolos embebidos. El código no impone una tarea ni una conexión nueva
por cada flujo NAPT, pero los límites internos de la pila siguen
condicionando la carga posible.

Los fallos se manejan según su importancia:
\begin{itemize}
  \item La inicialización imprescindible de NVS, red y listener DNS usa
  comprobaciones de error que pueden detener el arranque ante un fallo.
  \item La pérdida del módem activa recuperación por eventos y temporizador,
  conservando el AP para corregir la configuración.
  \item Un fallo al iniciar HTTP se registra; no hay actualmente un
  supervisor que reintente automáticamente el arranque de ese servidor.
  \item Un error del LED desactiva el indicador y deja continuar la red.
\end{itemize}

\section{Seguridad y alcance real de la aplicación}
\label{sec:seguridad}

La seguridad del enlace inalámbrico y la administración web son capas
distintas. El AP utiliza \code{WIFI_AUTH_WPA3_PSK} cuando se compila con
\code{CONFIG_ESP_WIFI_SOFTAP_SAE_SUPPORT}; en caso contrario utiliza
\code{WIFI_AUTH_WPA2_PSK}. En ambos casos el código requiere PMF. No es
un modo mixto automático WPA2/WPA3 para todos los clientes. La STA fija
un umbral mínimo WPA2 y \code{WPA3_SAE_PWE_BOTH}. Un cliente antiguo sin
la seguridad o PMF necesarios puede no asociarse.

El servidor usa HTTP en puerto80, sin inicio de sesión propio, HTTPS
ni protección CSRF implementada por la aplicación. No se enlaza
exclusivamente a la dirección AP; puede resultar accesible también por
la dirección STA desde la red del módem, según el acceso que permita
esa red. Poder alcanzar la interfaz HTTP permite enviar credenciales
nuevas; la contraseña del AP no es una autenticación administrativa
adicional para solicitudes que lleguen por otros caminos.

Los perfiles no habilitan cifrado de NVS ni de flash. Guardar una cadena
en NVS aporta persistencia, no confidencialidad frente a lectura física
de la memoria. El firmware evita imprimir contraseñas en sus mensajes
propios y la API no las publica, pero no incorpora un sistema de gestión
de secretos, borrado seguro de todos los búferes temporales o control de
acceso por usuario. La clave del AP inicial es conocida y está definida
en \path{components/wifi/wifi_internal.h}.

El listener DNS sí está enlazado a la IP del AP y filtra procedencias
de su /24. Estas comprobaciones limitan su alcance, pero no equivalen
a DNS autenticado. La reescritura de IDs sirve para correlacionar
transacciones; no convierte el protocolo en DNS cifrado.

Como futuras mejoras, podrían evaluarse autenticación administrativa,
restricción de interfaces HTTP, protección de solicitudes de cambio,
cifrado de almacenamiento con aprovisionamiento de claves y un formato
versionado de credenciales. También pueden estudiarse DNS por TCP y
failover de origen. Son propuestas, no funciones incluidas ni garantías
de esta versión. OTA tampoco está implementado.

\section{Verificación y diagnóstico práctico}
\label{sec:verificacion}

\subsection{Qué evidencia existe}

Las validaciones registradas en el proyecto tienen alcances diferentes:
\begin{itemize}
  \item \textbf{Compilación por target.} Se generaron configuraciones
  nuevas para ESP32 y ESP32-S3 con ESP-IDF 6.1 y se compilaron los dos
  perfiles. Esto comprueba dependencias, símbolos, selección de target
  y enlace de recursos. No equivale a programar ambas placas en cada revisión.
  \item \textbf{Parser en host.} Se compila el C de producción y se prueban
  límites, UTF-8, escapes, duplicados y entradas inválidas.
  \item \textbf{DNS en host.} El código de producción se adapta a sockets
  del host y se prueban siete escenarios con UDP real en loopback, puerto
  de prueba 1053: sin origen, preservación del mensaje, clientes con igual
  ID y respuestas invertidas, cambio de origen, respuesta desconocida,
  expiración y nueva pérdida de origen. El firmware usa el puerto 53.
  \item \textbf{Interfaz en host.} Se ejecuta el JavaScript real de la
  confirmación con DOM y peticiones simuladas: errores, reintentos, nuevas
  credenciales, recuperación y respuestas de estado inválidas.
  \item \textbf{Pruebas manuales reportadas en S3.} Se verificaron primer
  guardado, navegación desde el AP, SSID erróneo, contraseña errónea,
  recuperación con datos correctos y reinicio conservando NVS.
\end{itemize}

Los comandos se encuentran en \path{tests/README.md}. Las pruebas de host
no emulan el controlador de radio, DHCP de lwIP, NAPT real ni el
comportamiento de todos los teléfonos. No hay aquí una medición de
throughput, latencia, consumo, carga máxima con tres estaciones o
estabilidad durante días. No debe deducirse una cifra de rendimiento
del éxito de compilación ni de una navegación manual.

\subsection{Orden útil para localizar un problema}

\begin{enumerate}
  \item \textbf{Acceso local.} Comprobar asociación a \code{ESP32-NAT},
  dirección obtenida por DHCP y acceso explícito por HTTP a
  \code{192.168.2.1}. No se implementa una redirección universal de portal
  cautivo ni una captura general de todos los dominios.
  \item \textbf{Credenciales.} Confirmar que se recibió la página de
  guardado y revisar \code{connection_error}/\code{wifi_reason}. Un error
  201 orienta a red no encontrada; 15/204, al intercambio de claves.
  \item \textbf{Dirección de salida.} Revisar \code{connected} y el log de
  IPv4. Asociación sin IP sugiere revisar DHCP antes de atribuirlo a NAPT.
  \item \textbf{Redes diferentes.} Revisar \code{subnet_conflict}; corregir
  la superposición antes de probar DNS o navegación.
  \item \textbf{NAPT y resolución.} Revisar \code{nat_enabled} y
  \code{dns_ready}. Si una dirección externa es accesible pero los nombres
  no resuelven, estudiar el DNS de origen y sus restricciones.
  \item \textbf{Servicio del módem.} Si todo lo local está preparado,
  comprobar que la red del módem realmente tenga Internet. Un LED verde
  no descarta una falla del proveedor.
\end{enumerate}

Un \code{ECONNRESET} HTTP puede ser el cierre de una sesión por parte del
navegador; por sí solo no demuestra caída Wi-Fi. En cambio, las razones
de desconexión del controlador y los eventos IP permiten relacionar un
fallo de radio con el estado de salida. No se necesita borrar NVS como
paso rutinario de diagnóstico: hacerlo elimina la persistencia que se
pretende verificar.

\section{Cómo generar y mantener este documento}

Este archivo es autónomo. En Overleaf se crea un proyecto en blanco,
se sube \code{funcionamiento_tecnico.tex} y se selecciona como documento
principal, o se copia su contenido en \code{main.tex}. El compilador
debe ser \textbf{pdfLaTeX}. Diagramas, tablas y referencias están dentro
del mismo archivo. No se requieren imágenes externas, BibTeX,
\code{minted} ni ejecución de comandos externos. Una segunda compilación
actualiza índice y referencias cruzadas si la herramienta no la realiza
automáticamente.

Para mantenerlo alineado con el firmware, revisar especialmente
\path{components/wifi/wifi_internal.h}, \path{components/wifi/wifi.c},
\path{components/wifi/sta.c}, \path{components/wifi/ap.c},
\path{components/server/server.c}, \path{components/dns_proxy/dns_proxy.c},
\path{components/nvs/nvs.c}, \path{components/status_led/status_led.c},
\path{app/saved.html} y los tres archivos de defaults. El README raíz
es la guía de uso; este informe explica mecanismos y decisiones. Los
comentarios y nombres citados deben actualizarse cuando cambien los
contratos públicos, límites o condiciones de disponibilidad.

Las referencias siguientes describen las APIs y protocolos de soporte;
la autoridad para los valores y funciones específicas de la aplicación
es su código fuente. Las páginas de documentación \code{latest} pueden
evolucionar; las compilaciones documentadas aquí corresponden a ESP-IDF 6.1.

\begin{thebibliography}{9}
\bibitem{wifi} Espressif Systems.
\emph{ESP32-S3 Wi-Fi Driver: Overview, configuración y canal compartido APSTA.}
\url{https://docs.espressif.com/projects/esp-idf/en/latest/esp32s3/api-guides/wifi-driver/overview.html}

\bibitem{netif} Espressif Systems.
\emph{ESP-NETIF Programmer's Manual: direccionamiento, DHCP y NAPT.}
\url{https://docs.espressif.com/projects/esp-idf/en/latest/esp32s3/api-reference/network/esp_netif_programming.html}

\bibitem{nvs} Espressif Systems.
\emph{Non-Volatile Storage Library.}
\url{https://docs.espressif.com/projects/esp-idf/en/latest/esp32s3/api-reference/storage/nvs_flash.html}

\bibitem{rmt} Espressif Systems.
\emph{Remote Control Transceiver (RMT).}
\url{https://docs.espressif.com/projects/esp-idf/en/latest/esp32s3/api-reference/peripherals/rmt.html}

\bibitem{devkit} Espressif Systems.
\emph{ESP32-S3-DevKitC-1 v1.1 User Guide: diferencias de revisión y LED.}
\url{https://docs.espressif.com/projects/esp-dev-kits/en/latest/esp32s3/esp32-s3-devkitc-1/user_guide_v1.1.html}

\bibitem{build} Espressif Systems.
\emph{Build System: Target-dependent sdkconfig defaults.}
\url{https://docs.espressif.com/projects/esp-idf/en/latest/esp32/api-guides/build-system.html#target-dependent-sdkconfig-defaults}

\bibitem{dnsrfc} P. Mockapetris.
\emph{RFC1035: Domain Names---Implementation and Specification.}
Estructura de mensajes, identificadores, errores y truncamiento DNS.
\url{https://www.rfc-editor.org/rfc/rfc1035.html}

\bibitem{dnstcp} J. Dickinson et al.
\emph{RFC7766: DNS Transport over TCP---Implementation Requirements.}
Referencia para el transporte TCP que este proxy todavía no implementa.
\url{https://www.rfc-editor.org/rfc/rfc7766.html}
\end{thebibliography}

\end{document}
